% SEGMENT OLP-0615-S01
% Part: methods
% Chapter: induction
% Section: induction-on-N

\documentclass[../../../include/open-logic-section]{subfiles}

\begin{document}

\olfileid{mth}{ind}{inN}

\olsection{இயல் எண்கள்~$\Nat$ மீது தொகுத்தறிதல்}

மிக எளிய வடிவத்தில், எல்லா இயல் எண்களுக்கும்
பொருந்தும் முடிவுகளை நிறுவத் தொகுத்தறிதல் உதவுகிறது.
$0$ இலிருந்து தொடங்கி மீண்டும் மீண்டும்~$1$ ஐக்
கூட்டினால் இறுதியில் ஒவ்வோர் இயல் எண்ணையும் அடையலாம்
என்ற உண்மையை அது பயன்படுத்துகிறது. ஆகவே ஒவ்வோர்
எண்ணுக்கும் ஏதோ ஒன்று மெய் என்று நிறுவ, (1)~அது
$0$ க்கு மெய் என்றும், (2)~ஓர் எண்~$n$ க்கு
மெய்யாக இருக்கும் போதெல்லாம் அடுத்த எண்~$n+1$
க்கும் மெய் என்றும் காட்டலாம். ``எண்~$n$ க்கு
$P$ என்ற பண்பு உண்டு'' என்பதை $P(n)$ என்று
சுருக்கமாக எழுதுவோம் (``எண்~$k$ க்கு $P$
என்ற பண்பு உண்டு'' என்பதை $P(k)$ என்று
எழுதுவது போல). அப்போது எல்லா $n \in \Nat$
க்கும் $P(n)$ என்பதைத் தொகுத்தறிதலால் நிறுவுவது
பின்வருவனவற்றைக் கொண்டது:
\begin{enumerate}
\item $P(0)$ என்பதற்கான நிறுவல்;
\item தன்னிச்சையான~$k$ க்கு, $P(k)$ எனில்
      $P(k+1)$ என்பதற்கான நிறுவல்.
\end{enumerate}
இதைத் தெளிவாகப் பார்க்க (1),~(2) இரண்டும்
நம்மிடம் உள்ளன என்க. (1) மூலம் $P(0)$ மெய்.
(2) உம் இருந்தால், குறிப்பாக $P(0)$ எனில்
$P(0+1)$, அதாவது $P(1)$, என்று தெரியும்.
``தன்னிச்சையான~$k$ க்கு, $P(k)$ எனில்
$P(k+1)$'' என்ற பொதுக் கூற்றில்~$k$ க்குப்
பதிலாக~$0$ ஐ இட்டால் இது கிடைக்கும். ஆகவே
முன்னிலையை உறுதிப்படுத்தும் விதியால்~$P(1)$.
மீண்டும் (2) இல் இப்போது~$n$ க்கு~$1$
ஐக் கொண்டால், $P(1)$ எனில்~$P(2)$.
$P(1)$ ஐ ஏற்கெனவே நிறுவியிருப்பதால் அதே
விதியால்~$P(2)$ கிடைக்கிறது. இப்படியே தொடரலாம்.
எந்த எண்~$n$ க்கும் இதை $n$~முறை செய்தபின்
இறுதியில்~$P(n)$ ஐ அடைகிறோம். எனவே (1),~(2)
சேர்ந்து எல்லா $n \in \Nat$ க்கும்~$P(n)$
என்பதை நிறுவுகின்றன.

% SEGMENT OLP-0615-S02
ஓர் எடுத்துக்காட்டைப் பார்ப்போம். $n$ பகடைகளை
உருட்டும்போது கிடைக்கக்கூடிய வெவ்வேறு கூட்டுத்தொகைகள்
எத்தனை என்று அறிய விரும்புகிறோம் என்க. சற்றே
வினோதமாகத் தோன்றினாலும் $0$ பகடைகளிலிருந்து
தொடங்குவோம். பகடையே இல்லாவிட்டால், கிடைக்கக்கூடிய
கூட்டுத்தொகை ஒன்று மட்டுமே: புள்ளிகளே இல்லாததால்
கூட்டுத்தொகை~$0$. ஆகவே வெவ்வேறு முடிவுகளின்
எண்ணிக்கை~$1$. ஒரே ஒரு பகடை, அதாவது $n=1$,
இருந்தால் $1$ முதல்~$6$ வரை ஆறு முடிவுகள்.
இரு பகடைகளில் $2$ முதல் $12$ வரை எந்தக்
கூட்டுத்தொகையும் கிடைக்கும்; மொத்தம் $11$ முடிவுகள்.
மூன்று பகடைகளில் $3$ முதல் $18$ வரை எந்த
எண்ணும் கிடைக்கும்; அதாவது $16$ முடிவுகள்.
$1$, $6$, $11$, $16$: ஒரு முறை தெரிகிறது.
விடை $5n+1$ ஆக இருக்குமோ? $5n+1$ என்பது
சாத்தியமான அதிகபட்ச எண்ணிக்கையே: $n$ பகடைகளில்
கிடைக்கும் குறைந்த கூட்டுத்தொகை $n$ (எல்லாவற்றிலும்
$1$), மிகப் பெரியது $6n$ (எல்லாவற்றிலும்
$6$); இவ்விரண்டுக்கும் இடையில் $5n+1$
எண்களே உள்ளன.

\begin{thm}
  $n$ பகடைகளைக் கொண்டு $n$ முதல் $6n$
  வரையிலான சாத்தியமான $5n+1$ முடிவுகளையும்
  பெறலாம்.
\end{thm}

% SEGMENT OLP-0615-S03
\begin{proof}
$P(n)$ என்பது ``$n$ பகடைகளைக் கொண்டு
$n$ முதல்~$6n$ வரையிலான எந்த எண்ணையும்
பெற முடியும்'' என்ற கூற்று என்க. தொகுத்தறிதலைப்
பயன்படுத்த, பின்வருவனவற்றை நிறுவுவோம்:
\begin{enumerate}
\item \emph{தொகுத்தறிதலின் அடிப்படை நிலை} $P(1)$:
  ஒரே ஒரு பகடையில் $1$ முதல் $6$ வரை
  எந்த எண்ணையும் பெற முடியும்.
\item \emph{தொகுத்தறிதலின் அடுத்த படி}:
  எல்லா $k$ க்கும், $P(k)$ எனில்~$P(k+1)$.
\end{enumerate}

% SEGMENT OLP-0615-S04
(1) ஐ $6$ முகப் பகடையைப் பார்த்தே நிறுவலாம்.
அதற்கு ஆறு முகங்களும் உள்ளன; $1$ முதல்~$6$
வரையிலான ஒவ்வோர் எண்ணும் அதன் ஏதோ ஒரு முகத்தில்
காணப்படுகிறது. எனவே ஒரே பகடையைக் கொண்டு $1$
முதல்~$6$ வரை எந்த எண்ணையும் பெறலாம்.

(2) ஐ நிறுவ, நிபந்தனைக் கூற்றின் முற்பகுதியான
$P(k)$ ஐ எடுத்துக்கொள்கிறோம். இந்த எடுகோள்
\emph{தொகுத்தறிதல் எடுகோள்} எனப்படுகிறது.
இதைக் கொண்டு~$P(k+1)$ ஐ நிறுவ வேண்டும்.
$k+1$ பகடைகளின் சாத்தியமான முடிவுகளை,
$k$~பகடைகளின் முடிவுகளையும் கூடுதலாக உள்ள
$k+1$-ஆம் பகடையின் முடிவையும் கொண்டு
விளக்குவதுதான் கடினமான பகுதி. தொகுத்தறிதல்
எடுகோளைப் பயன்படுத்த வேண்டுமானால் இதைச்
செய்ய வேண்டும்.

% SEGMENT OLP-0615-S05
தொகுத்தறிதல் எடுகோளின்படி, $k$~பகடைகளைக்
கொண்டு $k$ முதல்~$6k$ வரை எந்த எண்ணையும்
பெறலாம். $(k+1)$-ஆம் பகடையில்~$1$
விழுந்தால் கூட்டுத்தொகையுடன் $1$ சேரும்.
ஆகவே $k$~பகடைகளை உருட்டி, பின்னர்
$(k+1)$-ஆம் பகடையில்~$1$ பெற்றால்,
$k+1$ முதல் $6k+1$ வரை எந்த மதிப்பையும்
பெறலாம். மீதமிருப்பவை என்ன? $6k+2$
முதல் $6k+6$ வரையிலான மதிப்புகள்.
முதல் $k$ பகடைகளிலும் $6$ பெற்று,
$(k+1)$-ஆம் பகடையில் $2$ முதல் $6$
வரையிலான ஓர் எண்ணைப் பெற்றால் இவற்றையும்
பெறலாம். இவ்விரண்டையும் சேர்த்தால் $k+1$
பகடைகளைக் கொண்டு $k+1$ முதல் $6(k+1)$
வரை எந்த எண்ணையும் பெறலாம். அதாவது,
தொகுத்தறிதல் எடுகோளான~$P(k)$ ஐப்
பயன்படுத்தி~$P(k+1)$ ஐ நிறுவியுள்ளோம்.
\end{proof}

% SEGMENT OLP-0615-S06
எண்கள், கணங்கள் முதலிய பொருள்களின் ஒரு தொடர்
தானே ``தொகுத்தறிதல் முறையில்'' வரையறுக்கப்பட்டிருந்தால்,
அதாவது $n$-ஆம் பொருளைக் கொண்டு $(n+1)$-ஆம்
பொருளை வரையறுத்திருந்தால், அதைப் பற்றிய கூற்றை
நிறுவத் தொகுத்தறிதலைப் பல நேரங்களில் பயன்படுத்துகிறோம்.
எடுத்துக்காட்டாக,~$n$ வரையிலான இயல் எண்களின்
கூட்டுத்தொகை~$s_n$ ஐப் பின்வருமாறு வரையறுக்கலாம்:
\begin{align*}
  s_0 & = 0\\
  s_{n+1} & = s_n + (n+1)
\end{align*}
இவ்வரையறையின்படி:
\begin{align*}
  s_0 & = 0,\\
  s_1 & = s_0 + 1 && = 1,\\
  s_2 & = s_1 + 2 && = 1 + 2 = 3\\
  s_3 & = s_2 + 3 && = 1 + 2 + 3 = 6, \text{ முதலியன.}
\end{align*}
இப்போது $s_n = n(n+1)/2$ என்பதைத் தொகுத்தறிதலால்
நிறுவலாம்.

% SEGMENT OLP-0615-S07
\begin{prop}
  $s_n = n(n+1)/2$.
\end{prop}

\begin{proof}
  (1) $s_0 = 0\cdot(0 + 1)/2$ என்பதையும்,
  (2) $s_k = k(k+1)/2$ எனில்
  $s_{k+1} = (k+1)(k+2)/2$ என்பதையும்
  நிறுவ வேண்டும். (1) தெளிவு. (2) க்கு
  $s_k = k(k+1)/2$ என்ற தொகுத்தறிதல்
  எடுகோளை எடுத்துக்கொள்கிறோம். அதைப் பயன்படுத்தி
  $s_{k+1} = (k+1)(k+2)/2$ என்று
  காட்ட வேண்டும்.

  $s_{k+1}$ என்ன? வரையறையின்படி
  $s_{k+1} = s_k + (k+1)$. தொகுத்தறிதல்
  எடுகோளின்படி $s_k = k(k+1)/2$. இதை
  முந்தைய சமன்பாட்டில் பதிலீடு செய்தபின் பின்னங்களின்
  எளிய கணக்கீடு மட்டும் மீதமிருக்கிறது:
  \begin{align*}
    s_{k+1} & = \frac{k(k+1)}{2} + (k+1) = {}\\
    & = \frac{k(k+1)}{2} + \frac{2(k+1)}{2} = {}\\
    & = \frac{k(k+1) + 2(k+1)}{2} = {}\\
    & = \frac{(k+2)(k+1)}{2}.
  \end{align*}
\end{proof}

% SEGMENT OLP-0615-S08
இங்கே முக்கியமான பாடம்: தொகுத்தறிதல் முறையில்
வரையறுக்கப்பட்ட $a_n$ போன்ற ஒரு தொடரைப்
பற்றிய கூற்றை நிறுவும்போது, தொகுத்தறிதல்
இயல்பான வழி. அப்படி வரையறுக்கப்படாதபோதும்
(பகடை முடிவுகளின் எடுத்துக்காட்டைப் போல),
~$k+1$ க்கான நிகழ்வை~$k$ க்கான
நிகழ்வுடன் ஏதோ ஒரு வகையில் தொடர்புபடுத்த
முடிந்தால் தொகுத்தறிதலைப் பயன்படுத்தலாம்.

\end{document}
